\documentclass[aspectratio=169]{beamer}
\usetheme{Madrid}
\usecolortheme{seahorse}
\setbeamertemplate{navigation symbols}{}
\setbeamertemplate{footline}[frame number]
\usepackage[T1]{fontenc}
\usepackage[utf8]{inputenc}
\usepackage{booktabs}
\usepackage{array}

\title{A Data-Management Perspective on Data Engineering Design Patterns}
\subtitle{Specifications, Semantics, and Execution-Based Evaluation\\ICDE 2027 tutorial proposal --- slide preview}
\author{Motaz K. H. Saad \and Chiara Rucco \and Antonella Longo}
\institute{DataLab, Dept.\ of Engineering for Innovation, University of Salento}
\date{Updated preview}

\begin{document}

% 1
\begin{frame}
  \titlepage
  \vspace{0.3em}
  \small 1.5 h or 3 h (3 h preferred, hands-on lab) \quad \textbullet \quad Two 90-minute slots in 3 h version
\end{frame}

% 2
\begin{frame}{Why design patterns for data engineering?}
  \begin{itemize}
    \item Teams repeatedly build pipelines: ingest $\rightarrow$ validate $\rightarrow$ transform $\rightarrow$ orchestrate $\rightarrow$ promote
    \item The discipline needs a shared, reusable vocabulary for these solutions
    \item Design patterns capture proven structure, trade-offs, and applicability
    \item \textbf{Common thread:} patterns as explicit \emph{specifications} (context, forces, obligations, consequences) that make LLM-generated, orchestrator-native artifacts checkable via execution-based scoring
  \end{itemize}
\end{frame}

% 3
\begin{frame}{Outline and 3-hour slot split}
  \small
  \begin{columns}[T]
    \begin{column}{0.5\textwidth}
      \textbf{Slot 1 (90 min)}
      \begin{itemize}
        \item P1 Foundations, related work, DM semantics (15)
        \item P2 Ingestion patterns (40)
        \item P3 ETLT/ELTL \& transformation (35)
      \end{itemize}
    \end{column}
    \begin{column}{0.5\textwidth}
      \textbf{Slot 2 (90 min)}
      \begin{itemize}
        \item P4 EADF / CI-CD promotion (15)
        \item P5 LLM pipelines \& evaluation (15)
        \item P6 Open challenges (5)
        \item P7 Hands-on lab (45)
        \item Debrief \& Q\&A (10)
      \end{itemize}
    \end{column}
  \end{columns}
  \vspace{0.4em}
  \footnotesize 1.5 h version: 90 min survey depth (P1--P6, no lab). Presenters: P1 Antonella (lead), Chiara, Motaz; P2 Chiara; P3--P4 Chiara (lead), Motaz (support); P5 Motaz (lead), Chiara (support); P6 Antonella (lead), all; P7 Chiara (lead), Motaz (support).
\end{frame}

% 4
\begin{frame}{Foundations and data-management semantics}
  \begin{itemize}
    \item POSA form: problem, context, forces, solution, consequences
    \item Pattern language: ingestion, transformation, \textbf{semantic layer}, \textbf{data quality}, orchestration, promotion, observability
    \item \textbf{Semantics covered across the tutorial:} delivery guarantees (exactly-once vs.\ at-least-once), idempotence and incremental loads, schema evolution, and lineage/provenance
    \item Positioned within broader literature: ETL/ELT foundations, lakehouse layering, Dataflow semantics, data quality/metadata/lineage systems, and LLM/agent benchmarks for data workflows
  \end{itemize}
\end{frame}

% 5
\begin{frame}{Ingestion patterns}
  \begin{itemize}
    \item \textbf{MIND} (metadata-driven): connectors declared as data; schema/metadata catalog; idempotent, resumable, incremental loads
    \item Cloud ingestion: object storage + queue + compute; batching vs.\ streaming; backpressure and retry
    \item Forces/semantics: exactly-once vs.\ at-least-once delivery, latency, cost, schema evolution, lineage/provenance, operational burden
    \item Case study: ingestion at scale in SESG (Integrated Platform for Enhanced Analysis of ESG Reports), an ICSC/IFAB-supported initiative requiring scalable, metadata-driven ingestion of heterogeneous document corpora
  \end{itemize}
\end{frame}

% 6
\begin{frame}{Transformation, semantic layer, and data quality patterns}
  \begin{itemize}
    \item ETL vs.\ ELT vs.\ ETLT vs.\ ELTL: stage structures, data layers, and trade-offs
    \item Enhanced variants (ETLT++/ELTL++): contracts, versioning, schema evolution, semantic-layer/semantic-curation obligations, lineage/provenance, continuous monitoring
    \item \textbf{Data-quality design patterns} as reusable specifications: profiling, rule-based validation, anomaly detection, issue triage, and governed remediation
    \item Positioned against lakehouse/medallion layering and foundational ETL/integration research
  \end{itemize}
\end{frame}

% 7
\begin{frame}{EADF: Environment-aware promotion in CI/CD}
  \begin{itemize}
    \item Promotion problem: environment-specific bindings as hidden technical debt
    \item EADF: four YAML artifacts (component, environment catalog, pipeline template, promotion policy) with late binding at deploy time
    \item Data-quality gates, schema evolution, lineage/provenance, stateful components
    \item Tool-neutral application (e.g. GitHub Actions, GitLab CI, Azure DevOps, Airflow, dbt)
  \end{itemize}
\end{frame}

% 8
\begin{frame}{LLM-generated pipelines and execution-based evaluation}
  \begin{itemize}
    \item Orchestrator landscape: Airflow, n8n, Dagster, Prefect, NiFi, Kestra
    \item Syntactic validity $\neq$ correctness: evaluate by executing artifacts and scoring sink-level data
    \item Execution-based benchmark: tasks, archived fixtures, one contract, containerized execution, repair rounds; failure taxonomy (artifact, execution, platform, harness)
    \item Pattern obligations as checks: idempotency, schema/contract conformance, quality gates
  \end{itemize}
\end{frame}

% 9
\begin{frame}{Hands-on lab (3 h, Slot 2)}
  \begin{itemize}
    \item \textbf{Lightweight orchestrators:} Prefect and Kestra (containerized)
    \item Flow: select task $\rightarrow$ generate orchestrator-native artifact $\rightarrow$ deploy/run $\rightarrow$ score sink records against reference output
    \item \textbf{Local path (optional):} Python 3.10+, Docker, Ollama; laptop $\geq$8\,GB RAM recommended
    \item \textbf{Fallback path:} hosted LLM endpoint (OpenAI-compatible) for generation; pre-generated artifacts and pre-computed scoring outputs, with optional instructor-hosted demo services
    \item Pre-pulled images/cached models; no-setup follow-along mode for all attendees
  \end{itemize}
\end{frame}

% 10
\begin{frame}{Open challenges and resources}
  \begin{itemize}
    \item Pattern mining, validation, and versioning
    \item Benchmark validity, contamination, and reproducibility
    \item Governance, provenance, cost, and human-in-the-loop pipelines
    \item Companion materials: benchmark harness/fixtures, ingestion/ETLT/EADF/LLM-generation repositories, data-quality design pattern materials (UNICC), slides and lab setup guide
  \end{itemize}
  \vspace{0.5em}
  \centering \url{https://github.com/motazsaad/data-eng-design-patterns-slides-lab}
\end{frame}

% 11
\begin{frame}
  \centering Thank you!
  \vspace{1em}

  \small Questions?

  \vspace{1em}
  \textit{Motaz K. H. Saad} --- \texttt{motazk.saad@unisalento.it}\\
  \textit{Chiara Rucco} --- \texttt{chiara.rucco@unisalento.it}\\
  \textit{Antonella Longo} --- \texttt{antonella.longo@unisalento.it}
\end{frame}

\end{document}
